\section{SYSTEM ARCHITECTURE AND WORKFLOW}
\label{sec:usage}

\subsection{Overview}
This edition of the automation runs \textbf{one Meso-NH configuration, once per night}: a single set of
nested domains (\texttt{NESTING} in \texttt{globals.var}, currently 3 -- outer domain plus two nested
levels), driven by one self-contained \texttt{bin/} tree and one \texttt{bin/conf.d/*.var} configuration
(Section~\ref{sec:config}). There is no built-in mechanism in this codebase for running several
configurations in sequence or for having one run's output conditionally replace another's; if a
deployment needs that (e.g.\ a coarser and a finer-resolution configuration chained together), it has to
be implemented as an external wrapper around this code, and documented separately.

\subsection{Manual and scheduled operation}
The whole procedure is a single entry point, \texttt{bin/run.sh}, invoked with no arguments:
\begin{lstlisting}
cd bin
./run.sh
\end{lstlisting}
For nightly unattended operation this is normally launched from cron, e.g.:
\begin{lstlisting}
0 20 * * * /path/to/bin/run.sh > /path/to/sim.out 2>&1
\end{lstlisting}
\texttt{run.sh} locates its own configuration via \texttt{conf.d/globals.var} relative to its own path,
so it must be run from within \texttt{bin/} (or by an absolute path that still resolves correctly). The
date being processed is controlled by \texttt{YEAR}/\texttt{MONTH}/\texttt{DAY} in \texttt{times.var}
(Section~\ref{sec:config}): under normal operation these are meant to track ``today'' (they default to
the system date at the top of the file), but \textbf{a hand-set override further down the file
currently takes precedence} -- see the note in Section~\ref{sec:config} on \texttt{times.var}. To
process a specific night (whether for the scheduled run or for manual reprocessing of a past date), edit
that override directly. There is no external process that substitutes today's date into the
configuration automatically; whatever \texttt{times.var} contains at the time \texttt{run.sh} sources it
is what gets used.

\subsection{Workflow (\texttt{run.sh})}
\texttt{bin/run.sh} performs, in order:
\begin{enumerate}
  \item \textbf{Startup}: kills any leftover Meso-NH-related processes from a previous, possibly stuck
        run (\texttt{pkill} on the executable names from \texttt{globals.var}), rotates the log
        directory, and sources \texttt{times.var} then \texttt{run.var}.
  \item \textbf{Initialization file retrieval}: sources \texttt{filedef.var} directly -- \textbf{this is
        where the actual wait happens}, in \texttt{run.sh}'s own process. \texttt{filedef.var}'s logic
        waits for the raw initialization files for the requested forcing hours to appear in
        \texttt{INIRAWDIR}, retries downloading them from a configurable backup server over SSH if they
        are late or missing, verifies each file's MD5 checksum, and merges the orography information
        from the analysis file into each forecast file with \texttt{grib\_copy} -- aborting with an
        alert mail if the files still have not arrived after a configurable timeout (\texttt{TIMELIMIT}
        in \texttt{times.var}). \texttt{run.sh} then sources \texttt{backup.var}.
  \item \textbf{Reference script generation}: only once the previous step has confirmed the
        initialization files are present, and only if \texttt{TESTRUN=0} and \texttt{PREPLISTTRUE=1}
        (and \texttt{NOPREP=0}\slash\texttt{NORUN=0}), \texttt{run.sh} invokes \texttt{bin/prep\_list.sh} as
        a sub-process. As explained in full in Section~\ref{sec:executables},
        \texttt{prep\_list.sh}'s own purpose is to generate a pair of self-contained, legacy-style
        reference scripts, not to fetch files; it re-sources \texttt{filedef.var} itself (in its own,
        separate shell) purely to know which files to write into those scripts, and since they are
        already present by this point, this second sourcing does not itself wait on anything.
        \texttt{run.sh} then moves the two generated scripts into \texttt{LOGDIR}.
  \item \textbf{PREP\_REAL phase}: for every nested domain, runs Meso-NH's \texttt{PREP\_REAL\_CASE}
        (and \texttt{SPAWNING} for nested domains) to build the initial and boundary conditions.
  \item \textbf{EXE phase}: runs the \texttt{MESONH} executable under \texttt{mpirun}, producing the
        raw \texttt{.lfi} model output for the whole forecast segment.
  \item \textbf{POSTPROC phase} (\texttt{bin/postproc.sh}, sourced): runs the \texttt{DIAG} step and
        \texttt{conv2dia}\slash\texttt{diaprog} to derive the astroclimatic parameters (optical
        turbulence $C_N^2$ and integrated parameters, wind, temperature, relative humidity, PWV, etc.)
        and to extract them as ASCII (FICVAL vertical profiles/cuts and 2D ground maps) at the points,
        domains and heights configured in \texttt{run.var}.
  \item \textbf{PLOT phase} (\texttt{bin/plot\_python.sh}, sourced): runs the Python plotting suite
        under \texttt{bin/conf.d/utils/Plot\_python/} (one script per parameter/plot type) to turn the
        ASCII output into the delivered figures.
  \item \textbf{Compression}: optionally \texttt{bzip2}-compresses the raw PREP\_REAL, RUN and DIAG
        \texttt{.lfi} outputs in parallel (GNU \texttt{parallel}), controlled by
        \texttt{DOZIPREAL}\slash\texttt{DOZIPOUT}\slash\texttt{DOZIPDIAG} in \texttt{globals.var}.
  \item \textbf{Move to staging}: moves the outputs selected for preservation
        (\texttt{BACKUP\_REAL}\slash\texttt{BACKUP\_RUN}\slash\texttt{BACKUP\_DIAG}\slash
        \texttt{BACKUP\_POSTPROC}\slash\texttt{BACKUP\_LOG} switches in \texttt{backup.var}) into a
        per-run staging directory (\texttt{PREPARE\_DIR}); everything else is deleted.
  \item \textbf{Upload/backup} (\texttt{bin/backup.sh}, sourced, unless deferred with
        \texttt{BACKUPLATER}): rsyncs the figures and ASCII ground outputs to the configured remote
        server(s) (Section~\ref{sec:config}, \texttt{backup.var}), and/or keeps a local copy, according
        to \texttt{backup.var}.
  \item \textbf{Notification and hand-off}: sends a completion e-mail with a per-phase timing summary;
        writes the day's configuration file for the short-timescale component into the \texttt{config/}
        subdirectory of \texttt{AUTOREGRESSION/}, and \texttt{rsync}s the whole \texttt{AUTOREGRESSION/}
        directory to the configured hand-off host(s), if that component is in use on this deployment
        (Section~\ref{sec:deps-install}). If so, the automation run must complete before it can operate
        correctly for that night.
\end{enumerate}
Figure~\ref{fig:runsh} summarizes this sequence.

\begin{figure}[H]
\centering
\begin{tikzpicture}[node distance=0.45cm]
\node[startstop] (start) {bin/run.sh start};
\node[process, below=of start] (filedef) {filedef.var\\wait \& retrieve init files};
\node[process, below=of filedef] (prep) {prep\_list.sh\\generate reference scripts\\(files already present)};
\node[mesonh, below=of prep] (real) {PREP\_REAL / SPAWNING (Meso-NH)};
\node[mesonh, below=of real] (exe) {EXE: MESONH (Meso-NH run)};
\node[process, below=of exe] (postproc) {postproc.sh\\DIAG + FICVAL extraction};
\node[process, below=of postproc] (plot) {plot\_python.sh\\plots};
\node[process, below=of plot] (compress) {compress (optional)\\bzip2, parallel};
\node[process, below=of compress] (stage) {stage outputs\\to PREPARE\_DIR};
\node[process, below=of stage] (backup) {backup.sh\\upload / local backup};
\node[startstop, below=of backup] (notify) {notify + hand-off};
\draw[wkarrow] (start) -- (filedef);
\draw[wkarrow] (filedef) -- (prep);
\draw[wkarrow] (prep) -- (real);
\draw[wkarrow] (real) -- (exe);
\draw[wkarrow] (exe) -- (postproc);
\draw[wkarrow] (postproc) -- (plot);
\draw[wkarrow] (plot) -- (compress);
\draw[wkarrow] (compress) -- (stage);
\draw[wkarrow] (stage) -- (backup);
\draw[wkarrow] (backup) -- (notify);
\end{tikzpicture}
\caption{Internal phase sequence of \texttt{bin/run.sh}. Grey boxes are Meso-NH itself; blue boxes are
automation scripts.}
\label{fig:runsh}
\end{figure}

Each of these phases can be individually disabled for debugging via the \texttt{NOPREP}\slash
\texttt{NORUN}\slash\texttt{NODIAG}\slash\texttt{NOFICVAL}\slash\texttt{NOPLOT} switches in
\texttt{globals.var}; this is intended for manual troubleshooting only.

\texttt{bin/pgdgen.sh} is a separate, standalone script (not called by \texttt{run.sh}) used to
(re-)generate the static PGD orography files for a site; it is normally run once, at installation time
(Section~\ref{sec:deps-install}), and only re-run if the domain geometry or the ground/topography input
files change.

The scripts \texttt{bin/Run\_lista.sh}, \texttt{bin/crossc2d.py} and \texttt{bin/multi.py} are
CPU/timing benchmarking helpers used to test the code against a list of past initialization dates; they
are not part of the nightly operational workflow.

\subsection{Directory layout}
The automation lives in one top-level working directory (\texttt{WORKDIR} in \texttt{globals.var}):
\begin{lstlisting}
bin/            : scripts (see Section 5)
bin/conf.d/     : configuration files, namelist templates, correction tables, compiled
                   fortran helpers (see Section 4)
PGD/            : static orography files, generated once by pgdgen.sh and persistent
                   across runs
\end{lstlisting}
The following sub-directories are created fresh by each run and removed/replaced at the end of it
(anything left behind after a failed run is there for post-mortem debugging only, and will be deleted
on the next successful run):
\begin{lstlisting}
REAL/     : PREP_REAL binary output and namelists
RUN/      : Meso-NH (EXE) binary output and namelists
DIAG/     : DIAG / conv2dia binary output and namelists
ASCII/    : ASCII output of the postprocessing (DAT, FICVAL, PLOT_DAT, AUTOREGRESSION
            sub-directories)
PLOT/     : figures produced by the Python plotting suite
LOG/      : logs of the current/last run
\end{lstlisting}
Outside \texttt{WORKDIR}, in the automation user's home directory, the following are used across runs:
\begin{lstlisting}
<init dir>/          : combined (orography + forecast) init files for the current night
                        (INIDIR)
<backup init cache>/ : local cache of files retrieved from the backup init server
                        (BKSERVER_LOCALPATH)
BACKUP_OUTPUT/        : per-night backup of ASCII/plot outputs (see backup.var)
BACKUP_FIGURES/        : local mirror of the delivered figures
<staging dir>/         : staging directory used while moving outputs (PREPARE_DIR)
AUTOREGRESSION/         : hand-off area shared with the short-timescale component,
                          if in use on this deployment
TOPOGRAPHY/             : source topography files used to build the PGD
<python virtualenv>/   : Python virtualenv used by the automation (Section 3)
\end{lstlisting}
The exact names in angle brackets above are whatever \texttt{globals.var} currently sets them to
(Section~\ref{sec:config}); they are not fixed by the code.

\subsection{Logging and alerting}
Each phase writes to its own timestamped log file under \texttt{LOGDIR} (main log, error log, and one
log per Meso-NH executable), configured in \texttt{globals.var}. A cumulative log recording start/end
time and duration of every phase of every run is appended to the path set by
\texttt{LOG\_TIME\_CUMULATIVE} in \texttt{globals.var}. Fatal problems abort the script and send an
alert e-mail (function \texttt{error}); non-fatal issues are logged and mailed without stopping the run
(function \texttt{error\_soft}).
